%!TEX TS-program = xelatex
\documentclass[10pt,a4paper]{scrartcl}  

% ---------- Definitely used ----------

%Fontspec - xetex standard
\usepackage{fontspec}

%proper language processung
\usepackage[ngerman]{babel}

%geometry, very usefull
\usepackage{geometry}

% ------ Might be of use somtimes -----

%math stuff
\usepackage{amsmath}

% ----------- What is this? -----------
\usepackage{xltxtra,xunicode}



% ---------- Global Formatting --------

% Clause format:
\setlength{\parindent}{0pt}
\setlength{\parskip}{6pt plus 2pt minus 1pt}

% Page format:
\pagestyle{plain}  

% table format:
% packages pandoc uses for tables
\usepackage{longtable, booktabs}

% aligns tables to the left
\setlength{\LTleft}{2em}

% Images
\usepackage{graphicx}
\makeatletter
\def\maxwidth{\ifdim\Gin@nat@width>\linewidth\linewidth\else\Gin@nat@width\fi}
\def\maxheight{\ifdim\Gin@nat@height>\textheight\textheight\else\Gin@nat@height\fi}
\makeatother
% Scale images if necessary, so that they will not overflow the page
% margins by default, and it is still possible to overwrite the defaults
% using explicit options in \includegraphics[width, height, ...]{}
\setkeys{Gin}{width=\maxwidth,height=\maxheight,keepaspectratio}
% Set default figure placement to htbp
\makeatletter
\def\fps@figure{htbp}
\makeatother


% make lists pretty
\providecommand{\tightlist}{%
\setlength{\itemsep}{0pt}\setlength{\parskip}{0pt}}


% Margins
\geometry{a4paper,top=3cm,bottom=3cm,left=3cm,right=3cm}

\usepackage[juratotoc]{scrjura}


% Custop section counters for protocol
% \setcounter{secnumdepth}{2}
% \counterwithout{subsection}{section}
% \renewcommand\thesection{\Roman{section}
% \renewcommand\theClause{title=§\arabic{subsection}
% \makeatletter
% \renewcommand{\autodot}{}% Remove all end-of-counter dots
% \makeatother

\usepackage{enumitem}
\setenumerate{label=\roman*. }

\renewcommand{\thesection}{\Roman{section}}

\RedeclareSectionCommand[
  runin=false,
  afterindent=false,
  beforeskip=2\baselineskip,
  afterskip=.3\baselineskip]{section}

% Heading styles:
% These commands keep the koma system from making stupid sans serif section headings
%\setkomafont{title}{\rmfamily\mdseries\upshape\normalsize}
%\setkomafont{sectioning}{\rmfamily\mdseries\upshape\normalsize}
%\setkomafont{descriptionlabel}{\rmfamily\mdseries\upshape\normalsize}

% fixing section title ovelapping with "TOP"
%\usepackage[titles]{tocloft}
%\addtolength{\cftsecnumwidth}{1cm}
%\renewcommand{\cftsecfont}{\bfseries \sffamily}

% TOC Heading
% \addto\captionsgerman{
%   \renewcommand{\contentsname}%
%     {Inhaltsverzeichnis}%
% }
% this could be considered changing
%\setkomafont{section}{\large}
%\setkomafont{subsection}{\normalsize}

% Moves floats to the top of the page (used for the signing field)
% Makes problems with tables
%\makeatletter
%\setlength{\@fptop}{0pt}
%\makeatother

% set up hyperref
\usepackage[breaklinks=true]{hyperref}
\hypersetup{colorlinks,%
citecolor=black,%
filecolor=black,%
linkcolor=black,%
urlcolor=blue}
\usepackage{url}

%The Title 
\title{\LARGE Geschäftsordnung des Studentischen Konvents der Universität Würzburg}
% emtying date so it does not show up
\date{}


\begin{document}

\maketitle

\begin{figure}[h]
	\includegraphics[width=\textwidth]{logo.png}
\end{figure}


\vspace{2cm}

{\centering \Large \sffamily in der Fassung vom 27.11.18 mit der letzen Änderung vom 02.10.2019}

\pagebreak

\setcounter{tocdepth}{4}
\tableofcontents

\pagebreak


\section{Rede-, Antrags-, Vorschlags- und
Stimmrecht}\label{rede--antrags--vorschlags--und-stimmrecht}

\begin{contract}

\Clause{title=Stimmrecht und
Stimmrechtsübertragungen}\label{stimmrecht-und-stimmrechtsuebertragungen}

Jedes Konventsmitglied hat Stimmrecht im Studentischen Konvent. Eine
schriftliche Stimmrechtsübertragung auf jedes beliebige andere
Konventsmitglied für einzelne Sitzungen oder Teile von Sitzungen ist
möglich. Jedes anwesende Konventsmitglied kann nur eine
Stimmrechtsübertragung annehmen. Stimmrechtsübertragungen sind der
vorsitzenden Person vor Beginn oder während der Sitzung in schriftlicher
Form und eigenhändig unterschrieben vorzulegen.

Alternativ können Stimmrechtsübertragungen als digitale Kopie des
unterschriebenen Dokuments bis zum Beginn der Sitzung per elektronischer Post
(E-Mail) der vorsitzenden Person übermittelt werden. Hierbei hat das Mitglied
dafür Sorge zu tragen, dass die Übertragung zugestellt wurde.


\Clause{title=Rede- Antrags- und
Vorschlagsrecht}\label{rede--antrags--und-vorschlagsrecht}

Alle Mitglieder des Studentischen Konvents haben Rede-, Antrags- und
Vorschlagsrecht und können Anträge zur Geschäftsordnung stellen. Alle
Studierenden der Julius-Maximilians Universität Würzburg haben Rede-,
Antrags- und Vorschlagsrecht. Anderen Personen kann das Rederecht auf
Antrag erteilt werden.

\end{contract}

\section{Wahlen}\label{wahlen}

\begin{contract}


\Clause{title=Allgemeines}\label{allgemeines}

Die Wahl der vorsitzenden Person und der Mitglieder der Sprecherinnen- und
Sprecherrates sowie dessen Vorsitz wird nach den §§45 -- 47 der Grundordnung der
Julius-Maximilians Universität durchgeführt.

Der Studentische Konvent kann den vertretenden Personen der Studierenden im
Senat Vorschläge für die Benennung von Kommissionsmitgliedern machen; die Wahl
der Vorschläge erfolgt nach §8.

Die Besetzung der Referate erfolgt ebenfalls nach §8.


\Clause{title=Rücktritt von Neuwahl}\label{ruxfccktritt-von-neuwahl}

Die vorsitzende Person oder ihre Stellvertretung und die sieben vom
Studentischen Konvent gewählten Mitglieder des Sprecherinnen- und
Sprecherrates sowie alle weiteren vom Studentischen Konvent gewählten
Personen können von ihrem Amt zurücktreten. Im Falle einer vorzeitigen
Ausscheidung aus dem Amt ist für den Rest der Amtszeit eine Nachwahl
binnen zweier Wochen durchzuführen.


\Clause{title=Konstruktives Misstrauensvotum}\label{konstruktives-misstrauensvotum}

  Der Studentische Konvent kann jeder von ihm gewählten Person das
  Misstrauen nur dadurch aussprechen, dass er mit der Mehrheit seiner
  Mitglieder eine*n Nachfolger*in für das zu besetzende Amt wählt. Einem
  Team kann nur in seiner Gesamtheit das Misstrauen ausgesprochen
  werden, als Nachfolger*in ist hier auch eine Einzelperson möglich.
  
	Der Antrag auf ein konstruktives Misstrauensvotum unter Nennung
  eines*r Kandidat*in oder mehrerer Kandidat*innen für die Nachfolge
  muss von mindestens einem Viertel der Mitglieder des Studentischen
  Konvents unterstützt werden und eine Woche vor der Sitzung der
  vorsitzenden Person vorliegen und von dieser in die Einladung
  aufgenommen werden. Ein Initiativantrag ist ausdrücklich
  ausgeschlossen.
     
	Der in Absatz (2) genannte Antrag kann mehrere Wahlvorschläge enthalten;
  der Studentische Konvent kann zudem während der Behandlung des Antrags
  mit je einem Viertel seiner Mitglieder weitere Kandidat*innen
  vorschlagen. In jedem Fall ist ein*e Kandidat*in gewählt, wenn sie
  oder er die Mehrheit der Mitglieder auf sich vereint. Es wird geheim
  gewählt.
     
  Anträge auf ein konstruktives Misstrauensvotum, die nicht den
  Voraussetzungen genügen, werden nicht auf die Tagesordnung gesetzt.
     
  Für die Abwahl der vorsitzenden Person des Studentischen Konvents und
  ihres*r Stellvertreter*in und der Mitglieder sowie der vorsitzenden
  Person des Sprecherinnen- und Sprecherrats gemäß Grundordnung gilt
  ebenfalls das Quorum sowie die obengenannten Voraussetzungen und
  Fristen für das Misstrauensvotum.


\Clause{title=Vorschlagsliste}\label{vorschlagsliste}

Die vorsitzende Person eröffnet die Liste der vorgeschlagenen Personen
und schließt sie, wenn keine weiteren Vorschläge vorliegen. Die Liste
der vorgeschlagenen Personen ist bis zur Abstimmung jederzeit auf
Verlangen von fünf Konventsmitgliedern erneut zu öffnen. Vor Beginn der
Abstimmung ist die Liste der vorgeschlagenen Personen zu verlesen.


\Clause{title=Personalbefragung und
Personaldebatte}\label{personalbefragung-und-personaldebatte}

Auf Verlangen eines Konventsmitglieds ist eine Personalbefragung
durchzuführen. Ebenso ist auf Verlangen eines Konventsmitglieds eine
Personaldebatte unter Ausschluss der Kandidat*innen durchzuführen. Die
Dauer von Personalbefragung und -debatte kann auf Antrag begrenzt
werden. Personaldebatten sind nichtöffentlich durchzuführen. Ein Antrag
auf Öffentlichkeit ist nicht möglich.


\Clause{title=Wahl von Referaten und Vorschlägen für Kommissionsmitglieder
}\label{wahl-von-referaten-und-kommissionsmitgliedern}

Gewählt werden: 
\begin{enumerate}
\item 
	Vorschläge für Mitglieder der
Kommissionen des Senats, der Universitätsleitung (UL) oder der
Erweiterten Universitätsleitung (EUL) sowie die Servicezentren
Professional School of Education (PSE) und Innovatives Lehren und
Studieren (ZiLS).
\item 
	Eine studentische Frauenbeauftragte und eine*n
Gleichstellungsbeauftragte*n. 
\item 
	Sowie Vorschläge für stellvertretende
Mitglieder bzw. Stellvertreter*innen für i. und ii.; es muss eine
stellvertretende studentische Frauenbeauftragte gewählt werden.

\end{enumerate}

Für
jedes Referat wird eine Person oder ein gemeinsam antretendes Team als
Referatsleitung gewählt. Ein solches Team wird als einzelner
Wahlvorschlag behandelt.

Es wird eine Vorschlagsliste nach §6
erstellt.

Gewählt ist bzw. sind entsprechend der Anzahl der zu
besetzenden Posten: 
\begin{enumerate}
\item Ein Wahlvorschlag für einen Vorschlag für
(stellvertretenende) Mitglieder, eine*n (stellvertretende*n)
studentische*n Beauftragte*n, wenn er die absolute Mehrheit der Stimmen
der Mitglieder des Studentischen Konvents auf sich vereint. In der
Stichwahl genügt die relative Mehrheit. 
\item Ein Wahlvorschlag für die
Besetzung eines Referates im ersten Wahlgang, wenn er die absolute
Mehrheit der Mitglieder des Studentischen Konvents auf sich vereint. In
der Stichwahl genügt die relative Mehrheit.
\item Bei Stimmengleichheit
findet eine Stichwahl statt; die Vorschlagsliste wird nicht neu
eröffnet. Bei einer erneuten Stimmengleichheit entscheidet das von der
die Wahl leitenden Person zu ziehende Los.
\item Umfasst die
Vorschlagsliste ausschließlich einen Vorschlag, beziehen sich i. und ii.
auf die Mehrheit der Ja-Stimmen gegen die Nein-Stimmen; anderenfalls
werden die Nein-Stimmen als ungültig gezählt. Leere Stimmzettel zählen
in jedem Fall als Enthaltung. Bei Stimmgleichheit gilt ein Vorschlag als
abgelehnt.
\end{enumerate}

Dieser Paragraph findet bei allen Wahlen, für die keine
gesonderte Regelung gilt, Anwendung.


\Clause{title=Der Sprecherinnen- und Sprecherrat
und Referate}\label{der-sprecherinnen--und-sprecherrat}

Alle Aufgaben des Sprecherinnen- und Sprecherrat sind Ressorts zugeteilt Diese
werden jeweils von einem Mitglied des Sprecherinnen- und Sprecherrats geleitet.
Ausgenommen sind die Studentischen Senator\*innen. Der Ressortzuschnitt wird
vom Studentischen Konvent per Beschluss festgelegt. Dieser Beschluss benötigt
eine Zwei-Drittel-Mehrheit.

Der Sprecherinnen- und Sprecherrat bestimmt die Besetzung der Ressortleitungen
selbst.

Ein Referat wird im studentischen Konvent beantragt und mit einer
Zwei-Drittel-Mehrheit beschlossen. Dabei werden die Aufgabengebiete des Referats
abschließend festgelegt. Die in der letzten Legislaturperiode bestehenden
Referate bleiben bestehen. Die Auflösung eines Referats kann nur durch eine
Zwei-Drittel-Mehrheit im Studentischen Konvent beschlossen werden.

Die Wahl der Referatsleitungen erfolgt nach und wird in der konstituierenden
Sitzung durchgeführt. Die zu wählenden Referatsleitungen sind vom Sprecherinnen-
und Sprecherrat in der ersten Juniwoche auszuschreiben. Wird für ein Referat
keine Leitung gewählt, übernimmt diese Aufgabe der Sprecherinnen- und
Sprecherrat. Wird die Einrichtung eines neuen Referats beschlossen, so ist die
Wahl der Referatsleitung für die unmittelbar nachfolgende Sitzung
hochschulöffentlich auszuschreiben.

Die Amtszeit der Referatsleitungen beginnt mit dem neuen
Wintersemester und beträgt ein Jahr.

Jedes Referat hat ein Mitglied des Sprecherinnen-und Sprecherrats
als Pat*in. Mindestens einmal im Semester nehmen diese Pat*innen an
einer Sitzung ihres Referats teil.
   
Sämtliche Referatsleitungen des Studentischen Konvents legen zum Ende
ihrer Amtszeit dem neu gewählten Studentischen Konvent einen
Rechenschaftsbericht vor. In diesem berichten sie über die Tätigkeit
des Referats des letzten akademischen Jahres. Der Bericht ist
schriftlich bis spätestens zwei Wochen vor der ersten ordentlichen
Sitzung einzureichen. Außerdem soll die Referatsleitung des letzten
akademischen Jahres in der ersten ordentlichen Sitzung des
Studentischen Konvents anwesend sein und für Rückfragen zur Verfügung
stehen. Den Referatsleitungen ist darüber hinaus Gelegenheit zu geben,
den Bericht persönlich vorzustellen. Hat die Referatsleitung im Lauf
des akademischen Jahres gewechselt, gilt dieser Absatz für die zum
jeweiligen Zeitpunkt amtierende Referatsleitung. Der Bericht hat
sachlich zu erfolgen und alle Veranstaltungen, Veröffentlichungen,
Aktionen usw. des jeweiligen Referats zu enthalten.

\end{contract}


\section{Gang der Verhandlung}\label{gang-der-verhandlung}

\begin{contract}


\Clause{title=Allgemeines}\label{allgemeines-1}

  Der Studentische Konvent ist von der vorsitzenden Person mindestens
  zweimal im Semester, nach Möglichkeit zu Beginn und gegen Ende der
  Vorlesungszeit, einzuberufen.

  Die vorsitzende Person beruft den Studentischen Konvent unter Angabe
  der Tagesordnung schriftlich ein. Die Mitglieder des Studentischen
  Konvents können auch mittels elektronischer Post (E-Mail) eingeladen
  werden. Die Frist der Einberufung beträgt mindestens eine Woche. In
  Fällen besonderer Dringlichkeit kann die vorsitzende Person die Frist
  abkürzen und auch mündlich oder fernmündlich einberufen. Die
  Mitglieder des Sprecherinnen- und Sprecherrats sowie der*die
  Protokollant*in sind gesondert zu laden. Die Einladung ist darüber
  hinaus unmittelbar hochschulöffentlich zur Verfügung zu stellen.

  Die vorsitzende Person des Studentischen Konvents ist verpflichtet,
  auf Verlangen von mehr als einem Viertel der Mitglieder des Gremiums
  innerhalb einer Frist von zwei Wochen zu einer Sitzung zu laden.
  Erfolgt diese Ladung nicht, ist jedes andere Konventsmitglied
  berechtigt, fristgerecht zu einer Sitzung zu laden. Nach Stattfinden
  dieser Sitzung liegt das Laderecht wieder bei der vorsitzenden Person.
  Eventuelle weitere nach Satz 2 erfolgte Ladungen sind rückwirkend
  ungültig.

  Die vorsitzende Person ist für die Ordnung verantwortlich und hat laut
  Versammlungsgesetz § 7 Hausrecht. Die vorsitzende Person kann die
  Person, die vom Verhandlungsgegenstand abschweift, mit Nennung des
  Namens zur Sache rufen. Ist die Person während einer Sitzung dreimal
  zur Sache gerufen und beim zweiten Male auf die Folgen eines dritten
  Rufes zur Sache hingewiesen worden, so ist ihr von der vorsitzenden
  Person das Wort zu entziehen. In der Debatte zum selben
  Verhandlungsgegenstand hat diese Person kein Rederecht mehr. Eine
  Person, die im Laufe einer Sitzung dreimal zur Ordnung gerufen wird,
  kann von der vorsitzenden Person des Sitzungsraumes verwiesen werden.

  Die Sitzungen des Studentischen Konvents sind grundsätzlich
  hochschulöffentlich. Ausgenommen hiervon sind Wahlen,
  Personalangelegenheiten sowie Tagesordnungspunkte, für die die
  Nichtöffentlichkeit per Antrag an die Geschäftsordnung beschlossen
  wurde. Die protokollierende Person und die Mitglieder des
  Sprecherinnen- und Sprecherrats sind grundsätzlich von der
  Nichtöffentlichkeit ausgenommen. Über weitere Ausnahmen beschließt der
  Studentische Konvent im Einzelfall.

  Es zählt zu den Pflichten der vorsitzenden Person sicherzustellen,
  dass für die Sitzung die folgenden Dinge bereitgestellt sind:

\begin{enumerate}
\item
  Stimmkarten mit Name des Mitglieds sowie Name der zugehörigen
  Gruppierung für alle Mitglieder
\item
  Stimmzettel, falls Wahlen für die entsprechende Sitzung ausgeschrieben
  wurden
\item
  ein Projektor
\item
  Verlängerungskabel und Mehrfachsteckdosen, um eine Stromversorgung für
  alle Mitglieder zu gewährleisten.
\end{enumerate}


\Clause{title=Tagesordnung}\label{tagesordnung}

  Die Tagesordnung muss folgende Punkte enthalten:
\begin{enumerate}[label=\alph*]
	\tightlist
\item Festlegung des*r
  Protokollant*in
\item Feststellung der Anwesenheit und der
  Beschlussfähigkeit
\item Genehmigung des Protokolls der vorangegangenen
  Sitzung
\item Genehmigung der Tagesordnung
\item Bericht aus dem
  Fachschaftenrat, den Kommissionen und Kollegialorganen sowie den
  Ausschüssen des Studentischen Konvents und Möglichkeit zur Diskussion
\item Bericht aus dem Sprecherinnen- und Sprecherrat und seinen Referaten
  und Möglichkeit zur Diskussion
\item Anträge
\item Verschiedenes
	\end{enumerate}

  Die Tagesordnung ist von den Mitgliedern des Studentischen Konvents
  mit einfacher Mehrheit zu genehmigen. Nach der Genehmigung der
  Tagesordnung muss die vorsitzende Person dann eine Umstellung der
  einzelnen Tagesordnungspunkte vornehmen, wenn diese Umstellung von
  einer Mehrheit der anwesenden Mitglieder des Studentischen Konventes
  genehmigt wird. Eine solche Umstellung kann per Antrag an die
  Geschäftsordnung beantragt werden.


\Clause{title=Beschlussfähigkeit}\label{beschlussfuxe4higkeit}



  Die Beschlussfähigkeit des Studentischen Konvents wird zu Beginn der
  Sitzung von der vorsitzenden Person festgestellt.

  Der Studentische Konvent ist beschlussfähig, wenn sämtliche Mitglieder
  ordnungsgemäß geladen sind und mehr als die Hälfte der Mitglieder
  physisch oder per Stimmrechtsübertragung anwesend ist. Andernfalls
  hebt die vorsitzende Person die Sitzung auf und beruft sie -- mit
  einwöchiger Ladungsfrist -- unter Beibehaltung der Tagesordnung binnen
  zweier Wochen neu ein; in diesem Fall ist der Studentische Konvent
  ohne Rücksicht auf die Zahl der anwesenden Stimmberechtigten
  beschlussfähig. Hierauf ist in der Einladung hinzuweisen. Schriftliche
  Stimmrechtsübertragungen werden bei der Feststellung von Anwesenheit
  und Stimmrecht von Mitgliedern mitberücksichtigt.

  Der Studentische Konvent gilt sodann, auch wenn sich die Zahl der
  anwesenden stimmberechtigten Mitglieder im Laufe der Sitzung
  verringert, als beschlussfähig, solange nicht ein Mitglied die
  Beschlussunfähigkeit geltend macht. Ist eine Angelegenheit wegen
  Beschlussunfähigkeit zurückgestellt worden und wird das Gremium zur
  Behandlung desselben Gegenstands erneut geladen, so ist es ohne
  Rücksicht auf die Zahl der anwesenden Stimmberechtigten
  beschlussfähig, wenn darauf in dieser Ladung hingewiesen worden ist.

  Stimmberechtigt sind die bei der Eröffnung der Abstimmung anwesenden
  stimmberechtigten Konventsmitglieder. Stimmrechtsübertragungen werden
  berücksichtigt.



\Clause{title=Dauer der Sitzungen}\label{dauer-der-sitzungen}




  Die Sitzungen des Studentischen Konvents sind spätestens um 1.00 Uhr
  zu beenden.

  Ein begonnener Tagesordnungspunkt wird abgeschlossen. Besteht ein
  Tagesordnungspunkt aus mehreren Unterpunkten, so gilt Satz 1 für jeden
  Unterpunkt.

  Der Studentische Konvent kann die Behandlung weiterer
  Tagesordnungspunkte bzw. deren Unterpunkte nur mit
  Zwei-Drittel-Mehrheit beschließen. Die der Abstimmung über die
  Behandlung weiterer Tagesordnungspunkte vorangehende Debatte ist auf 5
  Minuten zu beschränken.

  Zur Behandlung verbleibender Tagesordnungspunkte und deren Unterpunkte
	muss von der vorsitzenden Person gemäß §10 (2) der Geschäftsordnung eine
  Folgesitzung einberufen werden. Diese findet spätestens zwei Wochen
  nach der abgebrochenen Sitzung des Studentischen Konvents statt.

	Für die Tagesordnung der Folgesitzung gilt §11 (1) Es können keine
  Anträge gemäß §16 eingereicht werden.



\Clause{title=Leitung der Sitzung}\label{leitung-der-sitzung}

Mit Ausnahme der konstituierenden Sitzung eröffnet, leitet und
schließt die vorsitzende Person die Sitzungen des Studentischen
Konvents.

Die vorsitzende Person wird auf eigenen Wunsch oder bei Verhinderung
durch ihre*n Stellvertreter*in vertreten oder kann sich die Aufgaben mit ihr teilen.

Bei Abwesenheit der vorsitzenden Person und ihres*r Stellvertreter*in
kann ein beliebiges Mitglied des Studentischen Konvents die
Sitzungsleitung übernehmen. Erfolgt Widerspruch, so ist eine Leitung des
Studentischen Konvents für die aktuelle Sitzung mit einfacher Mehrheit
zu wählen. Wird niemand gewählt, so gilt die Sitzung als aufgehoben und
ist wie unter §12 (2) beschrieben neu zu laden.


\Clause{title=Reihenfolge der Redner*innen}\label{reihenfolge-der-rednerinnen}

Die vorsitzende Person führt eine Liste der Redner*innen und erteilt
das Wort in der Reihenfolge der Meldungen, wobei Redner*innen, die
sich zu diesem Verhandlungsgegenstand das erste Mal zu Wort melden,
vorgezogen werden.

Wenn auf eine direkte Frage eine direkte Antwort als sinnvoll
erscheint, kann die vorsitzende Person abweichend von der Redeliste
der gefragten Person unverzüglich eine einmalige Antwortmöglichkeit
geben.

Eine Zwischenfrage wird durch Kreuzen der Arme angezeigt. Die
vorsitzende Person fragt den*die Redner*in, ob er*sie die
Zwischenfrage zulässt.


\Clause{title=Fristgemäße Anträge}\label{fristgemuxe4uxdfe-antruxe4ge}

Anträge an den Studentischen Konvent sind fristgemäß, wenn sie drei
Werktage vor der Sitzung bis spätestens 17:00 Uhr schriftlich bei der
vorsitzenden Person des Studentischen Konvents eingereicht wurden. Der
Tag der Absendung der Anträge und der Tag der Sitzung werden nicht
mitgerechnet. Die vorsitzende Person leitet die Anträge noch am Tag des
Fristablaufs per elektronischer Post (E-Mail) in einem allgemein
gebräuchlichen Format an alle Mitglieder weiter. Außerdem stellt er*sie
alle eingehenden Anträge unmittelbar nach Fristablauf
hochschulöffentlich zur Verfügung. Anträge sind zu begründen, dies kann
auch mündlich zu Beginn der Behandlung des Antrags erfolgen.


\Clause{title=Initiativanträge}\label{initiativantruxe4ge}

Nach Antragsschluss können nur noch Initiativanträge in die Sitzung des
Studentischen Konvents eingebracht werden. Zum Zeitpunkt der Antragstellung muss
der Antragstext schriftlich vorliegen. Sie bedürfen der Unterstützung durch
mindestens fünf Mitglieder des Studentischen Konvents. Über ihre Behandlung
entscheidet der Studentische Konvent mit der einfachen Mehrheit der anwesenden
Mitglieder.


\Clause{title=Änderungsanträge}\label{uxe4nderungsantruxe4ge}

Änderungsanträge sind schriftlich bei der vorsitzenden Person des
Studentischen Konvents vor Beginn der Abstimmung über den Antrag
einzureichen. Auch sind Änderungsanträge von der antragstellenden Person
den Mitgliedern des Studentischen Konvents schriftlich zu unterbreiten.
Der*die Antragsteller*in des Antrags, zu dem Änderungsanträge vorliegen,
kann die Übernahme von Änderungsanträgen erklären. Vor der Abstimmung
über den Antrag muss über alle vorliegenden Änderungsanträge entschieden
werden. Liegen konkurrierende Änderungsanträge vor, so gilt §19 (4).


\Clause{title=Abstimmungen}\label{abstimmungen}

Der Studentische Konvent beschließt die Annahme von Anträgen per
Abstimmung.

Der Studentische Konvent beschließt mit der Mehrzahl der abgegebenen
Stimmen; Stimmenthaltung gelten als nicht abgegebene Stimmen. Bei
Stimmengleichheit gilt der Antrag als abgelehnt.

Vor Eröffnung der Abstimmung werden die zur Abstimmung stehenden
Anträge verlesen, soweit sie den Konventsmitgliedern nicht schriftlich
vorliegen.

Liegen mehrere Anträge zur selben Sache vor, so wird über den weitest
gehenden Antrag zuerst abgestimmt. Bei Finanzfragen ist über den am
wenigsten weitgehenden Antrag zuerst abzustimmen. Die Abstimmung
erfolgt grundsätzlich durch Handzeichen. Auf Antrag eines anwesenden
Konventsmitglieds muss geheim abgestimmt werden.

Während der Abstimmung kann niemand das Rederecht erhalten.

Unmittelbar im Anschluss an eine Abstimmung muss das
Abstimmungsergebnis auf Antrag von mindestens einem Viertel der
anwesenden Konventsmitglieder nochmals überprüft werden.
Gegebenenfalls ist die Abstimmung zu wiederholen.



\Clause{title=Protokoll}\label{protokoll}

Über die Sitzungen des Studentischen Konvents ist ein
Verlaufsprotokoll anzufertigen.


\begin{enumerate}
\item
  Der*die Protokollant*in wird jeweils in der ersten Sitzung des
  Studentischen Konvents im Semester für die Dauer eines Semesters
  gewählt.
\item
  Sofern sich keine Person gemäß der unter i. genannten Regelung findet
  oder der*die Protokollant*in verhindert ist, wird die protokollierende
  Person durch Losentscheid unter den anwesenden Mitgliedern oder durch
  freiwillige Meldung vor Eröffnung der Sitzung durch die vorsitzende
  Person bestimmt. Die Auswahl durch Los ist verpflichtend und kann nur
  unter Angabe eines gewichtigen Grunds aufgehoben werden.
\end{enumerate}

Von der Protokollpflicht befreit sind Mitglieder des Sprecherinnen-
und Sprecherrates, die vorsitzende Person und deren Stellvertreter*in
sowie weitere Mitglieder des Studentischen Konvents, die in der
laufenden Legislaturperiode bereits Protokoll geführt haben.

Der*die Protokollant*in soll spätestens vierzehn Tage nach der
Konventssitzung das Protokoll in sauber ausformulierter,
maschinenlesbarer Form der vorsitzenden Person zukommen lassen. Das
Protokoll wird, gegebenenfalls mit persönlichen Erklärungen dazu, mit
der Einladung zur nächsten Sitzung verschickt. Es wird zusätzlich auch
der vorsitzenden Person des Fachschaftenrats zugeschickt.

In die Protokolle sind die Beschlüsse wörtlich aufzunehmen.

Abstimmungsergebnisse sind im Protokoll zu vermerken. Bei deutlicher
Mehrheit ist eine genaue Auszählung nicht nötig, es sei denn, dies
wird von einem Mitglied verlangt.

Jedes in der Sitzung anwesende Konventsmitglied hat das Recht, seine
Auffassung zu die Sitzung betreffenden Angelegenheiten in einer
schriftlich begründeten persönlichen Erklärung darzulegen. Die
persönliche Erklärung ist dem Protokoll anzufügen.

Zu Beginn der nächsten Sitzung muss das Protokoll mit einfacher
Mehrheit genehmigt werden. Danach wird es von der vorsitzenden Person
innerhalb einer Woche hochschulöffentlich zur Verfügung gestellt. Wird
das Protokoll nicht genehmigt, erfolgt in der darauf folgenden Sitzung
erneut eine Abstimmung über eine gegebenenfalls überarbeitete Fassung.
Folgt in der laufenden Legislaturperiode keine Sitzung
mehr, so ist das Protokoll alleine zu verschicken. Über am Ende der
Legislaturperiode noch nicht beschlossene Protokolle beschließt der
Studentische Konvent der nächsten Legislaturperiode, dafür sind die
Protokolle dessen Mitgliedern mit der Sitzungseinladung zukommen zu
lassen.

\end{contract}


\section{Verfahren bei
Geschäftsordnungsanträgen}\label{verfahren-bei-geschuxe4ftsordnungsantruxe4gen}

\begin{contract}


\Clause{title=Äußerungen und Anträge zur
Geschäftsordnung}\label{uxe4uuxdferungen-und-antruxe4ge-zur-geschuxe4ftsordnung}

Äußerungen und Anträge zur Geschäftsordnung dürfen sich nur mit dem
Gang der Verhandlung befassen.

Äußerungen zur Geschäftsordnung sind insbesondere:
\begin{enumerate}
	\item ein Hinweis zur Geschäftsordnung,
\item eine Anfrage zur Geschäftsordnung sowie 
\item das Zurückziehen einer Anfrage oder eines Antrages.
	\end{enumerate}

Anträge zur Geschäftsordnung benötigen im Allgemeinen eine einfache
Mehrheit der abgegebenen Stimmen zur Annahme. Es handelt sich dabei
insbesondere um:


\begin{enumerate}

\item
  Antrag auf Vertagung. Seine Annahme hat zur Folge, dass der Punkt
  nicht weiter behandelt und auf die Tagesordnung der folgenden Sitzung
  gesetzt wird.
\item
  Antrag auf Nichtbefassung. Seine Annahme hat zur Folge, dass der
  (Unter-)Punkt nicht erörtert wird. Hierfür bedarf es einer
  Zwei-Drittel-Mehrheit der Anwesenden. Die Sitzung wird mit dem
  nächsten (Unter-)Punkt der Tagesordnung fortgesetzt.
\item
  Antrag auf Schließung der Debatte und sofortige Abstimmung.
\item
  Antrag auf Schließung der Redeliste. Nach Annahme des Antrags besteht
  die einmalige Möglichkeit, sich noch auf die Redeliste setzen zu
  lassen.
\item
  Antrag auf Beschränkung der Redezeit. Dieser Antrag muss unter Angabe
  der Maximalredezeit pro Redebeitrag gestellt werden. Die
Redezeit kann entweder für einen Verhandlungsgegenstand, einen
Tagesordnungspunkt oder die restliche Sitzung beschränkt werden 
\item
	Antrag auf Verhandlungspause. Der Antrag muss unter Angabe der
gewünschten Pausendauer gestellt werden. Die maximale Pausendauer
beträgt 15 Minuten. Während einer Sitzung können maximal 4 Pausen
beantragt werden. Längere oder weitere Pausen können mit einer
Zwei-Drittel-Mehrheit beschlossen werden. 
\item
	Antrag auf
Nichtöffentlichkeit. Dieser Antrag kann für die gesamte Sitzung,
einzelne Tagesordnungspunkte oder Verhandlungsgegenstände gestellt
werden.
\item
	Antrag auf Schließung der Sitzung. Seine Annahme hat zur Folge, dass die
	Sitzung sofort beendet wird und nach §13 Absätze (4) und (5) eine Folgesitzung
	einzuberufen ist. Hierfür bedarf es einer Zwei-Drittel-Mehrheit der
	Anwesenden.
\item
	Antrag auf geheime Abstimmung.
Dieser Antrag ist automatisch angenommen, eine Gegenrede ist nicht
möglich.
\end{enumerate}

Eine Wortmeldung zur Geschäftsordnung erfolgt durch Aufheben
beider Arme. Sie ist erst nach dem Ende des laufenden Redebeitrages,
dann aber unmittelbar zu behandeln. Mehrere Geschäftsordnungsanträge
werden in der Reihenfolge der Meldung behandelt. 

Die unter Abs. (3) iii.-v. genannten Anträge zur Geschäftsordnung dürfen nicht
direkt im Anschluss an einen eigenen Redebeitrag gestellt werden.

Ein Antrag zur Geschäftsordnung ist angenommen, wenn kein Mitglied des
Studentischen Konvents Widerrede einlegt. Die Widerrede kann formell oder
inhaltlich erfolgen, eine inhaltliche Widerrede hat dabei Vorrang. Die
inhaltliche Widerrede besteht aus einer kurzen Begründung, weshalb der Antrag
abgelehnt werden sollte.

\end{contract}


\section{Ausschüsse}\label{ausschuxfcsse}

\begin{contract}


\Clause{title=Wahl der Ausschüsse}\label{wahl-der-ausschuxfcsse}

Der Studentische Konvent kann zur Vorbereitung und zur Unterstützung
seiner Arbeit und für Untersuchungen ständige oder nichtständige
Ausschüsse einsetzen.

Den Ausschüssen dürfen nicht weniger als fünf Mitglieder angehören.
Sie müssen keine Mitglieder des Studentischen Konvents sein.

Die Mitglieder der Ausschüsse sowie deren Vorsitzende werden gewählt.
In den Ausschüssen sollten die verschiedenen Hochschulgruppen und
Fachschaften des Studentischen Konvents vertreten sein.

Treten gewählte Mitglieder aus einem Ausschuss zurück, so wählt der
Studentische Konvent in der nächsten Sitzung entsprechend viele neue
Mitglieder.

Auf Beschluss des Studentischen Konvents kann, um den Beitritt
weiterer Mitglieder zu ermöglichen, die Anzahl der Ausschussmitglieder
nachträglich erhöht oder bei Rücktritten auf nicht weniger als 
5 bzw. die Anzahl der noch aktiven Mitglieder verringert werden.



\Clause{title=Stimmrecht}\label{stimmrecht}

In den Ausschüssen haben alle Ausschussmitglieder Stimmrecht.


\Clause{title=Beschlussfassung der
Ausschüsse}\label{beschlussfassung-der-ausschuxfcsse}

Die Ausschüsse halten das Ergebnis ihrer Arbeit in Form von
schriftlichen Beschlüssen fest. Die Beschlüsse werden mit einfacher
Mehrheit gefasst. Das Ergebnis der Abstimmung ist anzufügen.


\Clause{title=Aufgaben und Rechte der
Ausschüsse}\label{aufgaben-und-rechte-der-ausschuxfcsse}

  Die dem Ausschuss vom Studentischen Konvent übertragenen Aufgaben sind
  gewissenhaft und ohne Verzögerung zu erledigen. Über ihre Erledigung
  ist dem Studentischen Konvent unverzüglich Bericht zu erstatten.

  Die Ausschüsse können darüber hinaus über jeden in ihren
  Arbeitsbereich fallenden Gegenstand verhandeln und Anträge im
  Studentischen Konvent einbringen.



\Clause{title=Sitzungen}\label{sitzungen}

  Die Ausschüsse werden von ihrem*r Vorsitzenden einberufen.

  Ein Ausschuss muss auf Verlangen von mindestens einem Viertel der
  Ausschussmitglieder einberufen werden. VI Informationsveranstaltung

\end{contract}

\section{Informationsveranstaltung}\label{informationsveranstaltung}

\begin{contract}

\Clause{title=Einladung zur
Informationsveranstaltung}\label{einladung-zur-informationsveranstaltung}

Die vorsitzende Person lädt einmal im Semester alle Studierenden zu
einer öffentlichen Informationsveranstaltung ein. Die Einladung ist
öffentlich bekanntzugeben. Die Mitglieder des Studentischen Konvents,
die gewählten Mitglieder des Sprecherinnen- und Sprecherrats, die
studentischen Mitglieder im Senat sowie die Mitglieder der
Fachschaftsvertretungen werden hierzu gesondert eingeladen. Die
Einladung der Studierenden erfolgt unter Angabe einer Tagesordnung durch
öffentlichen Aushang.


\Clause{title=Ablauf der
Informationsveranstaltung}\label{ablauf-der-informationsveranstaltung}




  Die vorsitzende Person des Studentischen Konvents leitet die
  Informationsveranstaltung oder trägt Sorge dafür, dass ein anderes
  Mitglied, bevorzugt ihre Stellvertreter*in, die Veranstaltung leitet.

  In der Informationsveranstaltung berichtet der Sprecherinnen- und
  Sprecherrat den Studierenden über seine Tätigkeiten. Die Studierenden
  haben die Möglichkeit, sich über die Arbeit ihrer Vertreter*innen in
  den Gremien über laufende Projekte, Referate und Veranstaltungen zu
  informieren und sich hierzu zu äußern.


\end{contract}


\section{Schlussbestimmungen}\label{schlussbestimmungen}

\begin{contract}


\Clause{title=Änderungen}\label{uxe4nderungen}

Änderungen der Geschäftsordnung bedürfen sowohl einer absoluten
Mehrheit aller Mitglieder des Konvents als auch eine
Zwei-Drittel-Mehrheit aller anwesenden Stimmen.

Anträge zur Änderung der Geschäftsordnung dürfen nicht als
Initiativanträge gestellt werden.


\Clause{title=Inkrafttreten}\label{inkrafttreten}

Diese Geschäftsordnung tritt am Tage ihrer Annahme durch den
Studentischen Konvent, am 27.11.2018, in Kraft.

\end{contract}


\end{document}
